%%%%%%%%%%%%%%%%%%%%%%%%%%%%%%%%%%%%%%%%%%%%%%%%%%%%%%%%%%%%
%% SDS 8102: Advanced Mathematical Statistics
%% Solutions: Laws of Large Numbers and Empirical Processes
%%%%%%%%%%%%%%%%%%%%%%%%%%%%%%%%%%%%%%%%%%%%%%%%%%%%%%%%%%%%

\documentclass[11pt]{article}

\usepackage[letterpaper,margin=1in]{geometry}
\usepackage{amsmath,amssymb,amsfonts}
\usepackage{mathtools}
\usepackage{xcolor}
\usepackage{enumitem}
\usepackage{fancyhdr}
\usepackage{titlesec}
\usepackage{tcolorbox}
\usepackage{microtype}
\usepackage[hidelinks]{hyperref}

\definecolor{MBZBlue}{HTML}{123B5D}
\definecolor{MBZTeal}{HTML}{00A6A6}
\definecolor{MBZLightBlue}{HTML}{EAF4F7}
\definecolor{MBZGold}{HTML}{D6A84B}
\definecolor{MBZGray}{HTML}{5F6B73}
\definecolor{MBZLightGray}{HTML}{F4F6F7}

\newcommand{\bbE}{\mathbb{E}}
\newcommand{\bbP}{\mathbb{P}}
\newcommand{\R}{\mathbb{R}}
\newcommand{\ind}{\mathbf{1}}
\DeclarePairedDelimiter{\abs}{\lvert}{\rvert}
\DeclarePairedDelimiter{\Norm}{\lVert}{\rVert}
\DeclareMathOperator{\Var}{Var}
\DeclareMathOperator{\Cov}{Cov}
\DeclareMathOperator{\card}{card}

\pagestyle{fancy}
\setlength{\headheight}{24pt}
\fancyhf{}
\fancyhead[L]{\textcolor{MBZBlue}{\textbf{SDS 8102}}}
\fancyhead[R]{\textcolor{MBZGray}{Solution Key}}
\fancyfoot[C]{\textcolor{MBZGray}{\thepage}}
\renewcommand{\headrulewidth}{0.5pt}
\renewcommand{\headrule}{%
  \hbox to\headwidth{\color{MBZTeal}\leaders\hrule height \headrulewidth\hfill}}

\titleformat{\section}{\Large\bfseries\color{MBZBlue}}{\thesection}{0.75em}{}
\titlespacing*{\section}{0pt}{1.5em}{0.8em}
\setlist[enumerate,1]{label=\textcolor{MBZBlue}{\textbf{\arabic*.}},leftmargin=*,itemsep=1em}
\setlist[enumerate,2]{label=\textcolor{MBZTeal}{\textbf{(\alph*)}},leftmargin=*,itemsep=0.6em}

\newtcolorbox{solutionbox}{
  colback=MBZLightBlue,colframe=MBZTeal,boxrule=0.8pt,arc=2mm,
  left=3mm,right=3mm,top=2mm,bottom=2mm}

\setlength{\parindent}{0pt}
\setlength{\parskip}{5pt}

\begin{document}

\begin{center}
{\color{MBZBlue}\rule{\textwidth}{1.2pt}}

\vspace{0.6em}
{\LARGE\bfseries\color{MBZBlue}SDS 8102: Advanced Mathematical Statistics}

\vspace{0.5em}
{\Large\bfseries\color{MBZTeal}Solution Key: Laws of Large Numbers \& Empirical Processes}

\vspace{0.4em}
{\large Fall 2026}

\vspace{0.2em}
{\large Mohamed bin Zayed University of Artificial Intelligence}

\vspace{0.6em}
{\color{MBZBlue}\rule{\textwidth}{1.2pt}}
\end{center}

\section{A Heavy-Tailed Moving-Average Process}

Let $m=\bbE[Y_t]$.

\begin{enumerate}
\item The Pareto density is
\[
f(y)=\alpha y^{-\alpha-1}\ind\{y\geq1\}.
\]
Thus
\[
\bbE[Y_t]
=\alpha\int_1^\infty y^{-\alpha}\,dy
=\frac{\alpha}{\alpha-1},
\]
whereas
\[
\bbE[Y_t^2]
=\alpha\int_1^\infty y^{1-\alpha}\,dy
=\infty
\]
because $\alpha<2$. Hence $\bbE[\varepsilon_t]=0$, and
\[
\bbE\abs{\varepsilon_t}
\leq \bbE[Y_t]+m<\infty.
\]
Moreover, if $Y_t\geq2m$, then $(Y_t-m)^2\geq Y_t^2/4$, so
\[
\bbE[\varepsilon_t^2]
\geq \frac14\bbE\left[Y_t^2\ind\{Y_t\geq2m\}\right]
=\infty.
\]

\item Linearity and the triangle inequality give
\[
\bbE[X_t]=0,
\qquad
\bbE\abs{X_t}
\leq(1+\abs{\rho})\bbE\abs{\varepsilon_0}<\infty.
\]
To verify the second moment carefully, choose $K>0$ such that
$\bbP(\abs{\rho\varepsilon_{t-1}}\leq K)>0$. By independence,
\[
\begin{aligned}
\bbE[X_t^2]
&\geq
\bbE\left[(\varepsilon_t+\rho\varepsilon_{t-1})^2
\ind\{\abs{\rho\varepsilon_{t-1}}\leq K,\ \abs{\varepsilon_t}>2K\}\right]\\
&\geq
\frac14\bbP(\abs{\rho\varepsilon_{t-1}}\leq K)
\bbE\left[\varepsilon_t^2\ind\{\abs{\varepsilon_t}>2K\}\right]
=\infty.
\end{aligned}
\]

\item $X_t$ is measurable with respect to
$\sigma(\varepsilon_t,\varepsilon_{t-1})$. If $\abs{t-s}>1$, the
two corresponding pairs of innovations are disjoint and hence independent.
Therefore $X_t$ and $X_s$ are independent.

\item Expanding and shifting indices,
\[
\begin{aligned}
\sum_{t=1}^nX_t
&=\sum_{t=1}^n\varepsilon_t
+\rho\sum_{t=1}^n\varepsilon_{t-1}\\
&=(1+\rho)\sum_{t=1}^{n-1}\varepsilon_t
+\varepsilon_n+\rho\varepsilon_0.
\end{aligned}
\]

\item Since the innovations are i.i.d., integrable, and mean zero, the SLLN gives
\[
\frac1n\sum_{t=1}^{n-1}\varepsilon_t\longrightarrow0
\qquad\text{a.s.}
\]
Writing $S_n=\sum_{t=1}^n\varepsilon_t$,
\[
\frac{\varepsilon_n}{n}
=\frac{S_n}{n}-\frac{n-1}{n}\frac{S_{n-1}}{n-1}\longrightarrow0
\qquad\text{a.s.}
\]
Also $\rho\varepsilon_0/n\to0$ a.s.

\item Dividing the identity in part 4 by (n) and applying part 5 yields
\[
\boxed{\frac1n\sum_{t=1}^nX_t\longrightarrow0\quad\text{a.s.}}
\]
Although $X_t$ has infinite variance, the proof only needs integrability of
the underlying i.i.d. innovations. A finite-variance version of the SLLN is
therefore not directly applicable to $X_t$.
\end{enumerate}

\section{The Classical Coupon-Collector Problem}

\begin{enumerate}
\item The time spent moving from (k-1) observed coupon types to (k)
types is $W_{k,n}=\tau_k-\tau_{k-1}$. The increments telescope:
\[
\sum_{k=1}^nW_{k,n}=\tau_n-\tau_0=T_n.
\]

\item After (k-1) distinct types have appeared, (n-k+1) types remain.
Thus
\[
p_{k,n}=\frac{n-k+1}{n},
\qquad
W_{k,n}\sim\operatorname{Geom}(p_{k,n}).
\]

\item By independence and the mean of a geometric variable,
\[
\bbE[T_n]
=\sum_{k=1}^n\frac1{p_{k,n}}
=n\sum_{j=1}^n\frac1j
=nH_n.
\]
Since $H_n/\log n\to1$,
\[
\frac{\bbE[T_n]}{n\log n}\longrightarrow1.
\]

\item With $j=n-k+1$,
\[
\begin{aligned}
\Var(T_n)
&=\sum_{k=1}^n\frac{1-p_{k,n}}{p_{k,n}^2}\\
&=\sum_{j=1}^n\frac{1-j/n}{(j/n)^2}
\leq n^2\sum_{j=1}^n\frac1{j^2}
\leq\frac{\pi^2}{6}n^2.
\end{aligned}
\]

\item Consequently,
\[
\Var\left(\frac{T_n}{n\log n}\right)
\leq\frac{\pi^2}{6\log^2n}\longrightarrow0.
\]
For every $\varepsilon>0$, Chebyshev's inequality gives
\[
\bbP\left(
\abs{\frac{T_n-\bbE[T_n]}{n\log n}}>\varepsilon
\right)
\leq\frac{\pi^2}{6\varepsilon^2\log^2n}\longrightarrow0.
\]

\item Decompose
\[
\frac{T_n}{n\log n}
=\frac{T_n-\bbE[T_n]}{n\log n}
+\frac{\bbE[T_n]}{n\log n}.
\]
The first term converges to zero in probability and the second to one.
Therefore
\[
\boxed{\frac{T_n}{n\log n}\longrightarrow1\quad\text{in probability}.}
\]
Thus the collection time is asymptotically $n\log n$ on its natural
first-order scale.
\end{enumerate}

\section{A Two-Series Theorem for Positive Random Variables}

Fix $c>0$ and put $W_n=V_n\ind\{V_n\leq c\}$. Kolmogorov's
three-series theorem says that $\sum_nV_n$ converges a.s. if and only if
\[
\sum_n\bbP(V_n>c)<\infty,
\qquad
\sum_n\bbE[W_n]\ \text{converges},
\qquad
\sum_n\Var(W_n)<\infty.
\]
Because $W_n\geq0$, convergence of the second series is exactly
$\sum_n\bbE[W_n]<\infty$. Furthermore,
\[
\Var(W_n)\leq\bbE[W_n^2]\leq c\,\bbE[W_n].
\]
Hence the variance series follows automatically from the expectation series.
The three-series criterion therefore reduces to
\[
\boxed{
\sum_nV_n<\infty\ \text{a.s.}
\iff
\sum_n\bbP(V_n>c)<\infty
\ \text{and}\ 
\sum_n\bbE[V_n\ind\{V_n\leq c\}]<\infty.}
\]

\section{Equivalent Criteria for Summability}

For every $x\geq0$,
\[
\frac12\min\{x,1\}
\leq\frac{x}{1+x}
\leq\min\{x,1\}.
\]
Indeed, consider separately $0\leq x\leq1$ and $x>1$. Taking
expectations and summing shows immediately that
\[
\sum_n\bbE\left[\frac{X_n}{1+X_n}\right]<\infty
\iff
\sum_n\bbE[\min\{X_n,1\}]<\infty.
\]

If the latter holds, Tonelli's theorem gives
\[
\bbE\left[\sum_n\min\{X_n,1\}\right]<\infty,
\]
so $\sum_n\min\{X_n,1\}<\infty$ a.s. In particular, $X_n>1$
occurs only finitely often, and hence
\[
\sum_nX_n<\infty\qquad\text{a.s.}
\]
Conversely, if $\sum_nX_n<\infty$ a.s., the two-series result with
$c=1$ yields
\[
\sum_n\bbE[X_n\ind\{X_n\leq1\}]<\infty,
\qquad
\sum_n\bbP(X_n>1)<\infty.
\]
Since
\[
\bbE[\min\{X_n,1\}]
=\bbE[X_n\ind\{X_n\leq1\}]+\bbP(X_n>1),
\]
condition (iii) follows. Thus (i), (ii), and (iii) are equivalent.

\section{The Strong Law and \(L^1\) Convergence}

Let $\mu=\bbE[X_1]$. Since $X_1\in L^1$, the de la
Vall\'ee--Poussin criterion provides an increasing convex function
$\Phi:[0,\infty)\to[0,\infty)$ such that
\[
\frac{\Phi(t)}{t}\longrightarrow\infty,
\qquad
\bbE\Phi(\abs{X_1})<\infty.
\]
By convexity and Jensen's inequality,
\[
\Phi\left(\abs{\frac{S_n}{n}}\right)
\leq
\Phi\left(\frac1n\sum_{i=1}^n\abs{X_i}\right)
\leq
\frac1n\sum_{i=1}^n\Phi(\abs{X_i}).
\]
Therefore
\[
\sup_n\bbE\Phi\left(\abs{\frac{S_n}{n}}\right)
\leq\bbE\Phi(\abs{X_1})<\infty.
\]
Another application of the de la Vall\'ee--Poussin criterion shows that
$\{S_n/n:n\geq1\}$ is uniformly integrable. The SLLN gives
$S_n/n\to\mu$ a.s., hence in probability. Uniform integrability then
upgrades this convergence to $L^1$:
\[
\boxed{
\bbE\abs{\frac{S_n}{n}-\mu}\longrightarrow0.}
\]

\section{The Empirical Distribution Function at a Fixed Point}

\begin{enumerate}
\item Put $I_i(x)=\ind\{X_i\leq x\}$. Then the $I_i(x)$ are i.i.d.
Bernoulli$(F(x))$, so
\[
\bbE[\widehat F_n(x)]=F(x),
\qquad
\Var(\widehat F_n(x))
=\frac{F(x)(1-F(x))}{n}.
\]

\item For each fixed $x$, the SLLN applied to $I_i(x)$ gives
\[
\widehat F_n(x)=\frac1n\sum_{i=1}^nI_i(x)\longrightarrow F(x)
\qquad\text{a.s.}
\]

\item Independence across sample indices implies
\[
\Cov(\widehat F_n(x),\widehat F_n(y))
=\frac1n\Cov(I_1(x),I_1(y)).
\]
Since $I_1(x)I_1(y)=I_1(x\wedge y)$,
\[
\boxed{
\Cov(\widehat F_n(x),\widehat F_n(y))
=\frac{F(x\wedge y)-F(x)F(y)}{n}.}
\]

\item Pointwise convergence controls each predetermined $x$, while the
supremum ranges over an uncountable set and its maximizing point may depend
on the sample and on $n$. Thus pointwise convergence alone does not imply
uniform convergence; an additional argument exploiting monotonicity or
empirical-process bounds is required.
\end{enumerate}

\section{A Uniform Law over a Finite Class}

\begin{enumerate}
\item For fixed $f$, the variables $f(X_i)$ are independent and lie in
$[0,1]$. Hoeffding's inequality gives
\[
\bbP(\abs{P_nf-Pf}>\varepsilon)
\leq2e^{-2n\varepsilon^2}.
\]

\item By the union bound,
\[
\begin{aligned}
\bbP\left(\sup_{f\in\mathcal F}\abs{P_nf-Pf}>\varepsilon\right)
&\leq\sum_{j=1}^M
\bbP(\abs{P_nf_j-Pf_j}>\varepsilon)\\
&\leq2Me^{-2n\varepsilon^2}.
\end{aligned}
\]

\item Set the last display equal to $\delta$. Solving for
$\varepsilon$ gives, with probability at least $1-\delta$,
\[
\boxed{
\sup_{f\in\mathcal F}\abs{P_nf-Pf}
\leq\sqrt{\frac{\log(2M/\delta)}{2n}}.}
\]

\item For each fixed $\varepsilon>0$,
\[
\bbP\left(\sup_{f\in\mathcal F_n}\abs{P_nf-Pf}>\varepsilon\right)
\leq2\exp\{\log M_n-2n\varepsilon^2\}.
\]
This tends to zero whenever
\[
\boxed{\log M_n=o(n).}
\]
Thus this condition guarantees convergence in probability.
\end{enumerate}

\section{Rademacher Complexity and Uniform Convergence}

Write $P_n'f=n^{-1}\sum_{i=1}^nf(X_i')$.

\begin{enumerate}
\item Since $Pf=\bbE[P_n'f\mid X_1,\ldots,X_n]$, conditional Jensen's
inequality gives
\[
\begin{aligned}
\sup_{f\in\mathcal F}\abs{P_nf-Pf}
&\leq
\bbE\left[
\sup_{f\in\mathcal F}\abs{P_nf-P_n'f}
\,\middle|\,X_1,\ldots,X_n
\right].
\end{aligned}
\]
Taking expectation proves the claim.

\item Put $\mathcal F_{\pm}=\mathcal F\cup(-\mathcal F)$. Then
\[
\sup_{f\in\mathcal F}\abs{(P_n-P_n')f}
=\sup_{f\in\mathcal F_{\pm}}(P_n-P_n')f.
\]
The joint distribution is unchanged when each pair $(X_i,X_i')$ is
swapped according to an independent Rademacher sign. Hence
\[
\begin{aligned}
\bbE\sup_{f\in\mathcal F_{\pm}}(P_n-P_n')f
&=\bbE\sup_{f\in\mathcal F_{\pm}}
\frac1n\sum_{i=1}^n\varepsilon_i(f(X_i)-f(X_i'))\\
&\leq
\bbE\sup_{f\in\mathcal F_{\pm}}\frac1n\sum_{i=1}^n\varepsilon_if(X_i)
+\bbE\sup_{f\in\mathcal F_{\pm}}\frac1n\sum_{i=1}^n(-\varepsilon_i)f(X_i')\\
&=2\bbE[\widehat{\mathfrak R}_n(\mathcal F_{\pm})].
\end{aligned}
\]

\item Conditional on the data, independence of the signs and
$\cosh u\leq e^{u^2/2}$ give
\[
\begin{aligned}
\bbE_\varepsilon\exp\left\{
\lambda\sum_{i=1}^n\varepsilon_if_j(X_i)\right\}
&=\prod_{i=1}^n\cosh(\lambda f_j(X_i))\\
&\leq\exp\left\{\frac{\lambda^2}{2}
\sum_{i=1}^nf_j(X_i)^2\right\}.
\end{aligned}
\]

\item Let
$R=\max_j(\sum_i f_j(X_i)^2)^{1/2}$. For every $\lambda>0$,
\[
\begin{aligned}
\bbE_\varepsilon\max_{1\leq j\leq M}
\sum_{i=1}^n\varepsilon_if_j(X_i)
&\leq\frac1\lambda
\log\sum_{j=1}^M
\bbE_\varepsilon e^{\lambda\sum_i\varepsilon_if_j(X_i)}\\
&\leq\frac{\log M}{\lambda}+\frac{\lambda R^2}{2}.
\end{aligned}
\]
Optimizing at $\lambda=\sqrt{2\log M}/R$ and dividing by $n$ yields
\[
\widehat{\mathfrak R}_n(\mathcal F)
\leq\frac{R\sqrt{2\log M}}{n}.
\]
Since $0\leq f_j\leq1$, $R\leq\sqrt n$, and therefore
\[
\widehat{\mathfrak R}_n(\mathcal F)
\leq\sqrt{\frac{2\log M}{n}}.
\]

\item The class $\mathcal F_{\pm}$ has at most $2M$ elements and all
its functions have absolute value at most one. Parts 2 and 4 give
\[
\boxed{
\bbE\sup_{f\in\mathcal F}\abs{P_nf-Pf}
\leq2\sqrt{\frac{2\log(2M)}{n}}.}
\]

\item Define
\[
Z(X_1,\ldots,X_n)=\sup_{f\in\mathcal F}\abs{P_nf-Pf}.
\]
Replacing one observation changes every $P_nf$ by at most $1/n$, and
hence changes (Z) by at most (1/n). McDiarmid's inequality gives
\[
\bbP(Z-\bbE Z>t)\leq e^{-2nt^2}.
\]
Taking $t=\sqrt{\log(1/\delta)/(2n)}$ and using part 5 proves that, with
probability at least $1-\delta$,
\[
\boxed{
Z\leq
2\sqrt{\frac{2\log(2M)}{n}}
+\sqrt{\frac{\log(1/\delta)}{2n}}.}
\]
Rademacher complexity measures how well a class can correlate with random
signs on the observed sample. A richer class can fit these signs more
closely and has larger complexity; a smaller complexity therefore yields a
tighter uniform deviation bound.
\end{enumerate}

\end{document}
